\documentclass[a4paper,12pt]{report}
\input{tasks_preamble_en.tex}
\setcounter{secnumdepth}{0}
\begin{document}

\begin{center}
{\Large\bfseries Learning objectives for the lesson \textup{3}}\\[2mm]
{\large Simplicial complexes, homology, Cech and Vietoris-Rips complexes}\\[2mm]
{\small Minicourse <<Gromov--Hausdorff distance: applications and lower bounds>>, lesson on presentation \textup{3} (chapter \textup{3} of the minicourse)}
\end{center}

\section*{Lesson plan (\textup{50} minutes)}

\begin{itemize}[leftmargin=*]
\item Problem \textup{1} (cycles, boundaries and minor homologies) --- \textup{8} minutes\textup{;}
\item Problem \textup{2} (nerve lemma) --- \textup{10} minutes\textup{;}
\item Problem \textup{3} (close simplicial mappings) --- \textup{8} minutes\textup{;}
\item Problem \textup{4} (stability of homology of Vietoris--Rips complexes) --- \textup{12} minutes\textup{;}
\item Problem \textup{5} (paracompactness and partitioning of unity) --- \textup{6} minutes\textup{;}
\item Problem \textup{6} (higher homology of connected manifolds) --- \textup{6} minutes.
\end{itemize}

At the presentation \textup{3} definitions were given and key results of the chapter \textup{3} were formulated. The inclusions $\VR(X,r) \subset \CechAmb{r}{X}{Z} \subset \VR(X,2r)$ and the key lemma on close spaces were proved. The homotopy invariance of homology, the nerve lemma, paracompactness of metric spaces, close simplicial maps and higher homology of connected manifolds are accepted without proof. Tasks \textup{1}--\textup{6} close exactly these gaps.

\section*{Reference information}

\begin{itemize}[leftmargin=*]
\item \emph{The chain group} $C_k(\Sigma; G)$ of an abstract simplicial complex $\Sigma$ with coefficients in the Abelian group $G$ is generated by the symbols $g[v_0, \dots, v_k]$ with the relations
\begin{gather*}
g_1[v_0,\dots,v_k] + g_2[v_0,\dots,v_k] = (g_1+g_2)[v_0,\dots,v_k], \\
g[v_{\pi(0)},\dots,v_{\pi(k)}] = \sgn(\pi)\, g[v_0,\dots,v_k].
\end{gather*}
\item \emph{Boundary operator} \textup{:}
\[
\partial_k\bigl(g[v_0,\dots,v_k]\bigr) = \sum_{i=0}^k (-1)^i\, g[v_0,\dots,\widehat{v_i},\dots,v_k], \qquad g \in G,
\]
and \emph{homology} is $H_k(\Sigma;G) = \ker \partial_k / \operatorname{im} \partial_{k+1}$.
\item \emph{Nerve} $\Nerve{\mathcal U}$ of the family $\mathcal U = \{U_\alpha\}$ \textup{:} simplexes are finite subfamilies with non-empty intersection\textup{;} a covering \emph{good} if all finite non-empty intersections of its elements are contractible.
\item \emph{External Cech complex} $\CechAmb{r}{X}{Z}$ \textup{:} simplices --- finite subsets of $\sigma \subset X$ for which there is $z \in Z$ with $|zx|_Z < r$ for all $x \in \sigma$ \textup{(}that is, this is the nerve of the family of balls $\{B(x;r)\}_{x \in X}$ in $Z$\textup{)}.
\item \emph{Vietoris--Rips complex} $\VR(X,r)$ \textup{:} simplices are finite subsets of diameter $< r$.
\item Two simplicial maps are \emph{close} \textup{(}contiguous\textup{)} if for each simplex $\sigma$ the set $f(\sigma) \cup g(\sigma)$ lies in some simplex.
\item \emph{Paracompactness}\textup{:} a space is paracompact if each of its open coverings has a locally finite open refinement \textup{(}refinement is a covering, each element of which lies in some element of the original\textup{)}. Every metric space is paracompact \textup{(}Stone's theorem\textup{)}.
\item \emph{Higher homology of manifolds}\textup{:} for a connected $n$-dimensional manifold $M$ without boundary $H_n(M;\Z_2) = \Z_2$ for compact $M$ and $H_n(M;\Z_2) = 0$ for non-compact $M$ \textup{(}Hatcher, Theorem \textup{3.26} and Proposition \textup{3.29}\textup{)}.
\end{itemize}

\section*{Tasks}

\begin{task}[cycles, boundaries and minor homologies]
Let $\Sigma$ be an abstract simplicial complex, $G$ an Abelian group, $C_k(\Sigma; G)$ a group of $k$-dimensional chains, and $\partial_k$ a boundary operator.
\begin{enumerate}[label=\textup{(\alph*)}]
\item Check that $\partial_{k-1} \circ \partial_k = 0$ for all $k \ge 2$, and deduce from this that the \emph{cycles} subgroup $Z_k = \ker \partial_k$ contains the \emph{boundaries} subgroup $B_k = \operatorname{im} \partial_{k+1}$, so the quotient group $H_k(\Sigma; G) = Z_k / B_k$ is defined. Explain why the definition of circuits requires the coefficient $g \in G$ for the symbol $[v_0, \dots, v_k]$.
\item Prove that for a connected $\Sigma$ the completion homomorphism $\varepsilon\bigl(g[v]\bigr) = g$ passes through $H_0$ and defines an isomorphism $H_0(\Sigma; G) \cong G$\textup{;} for a complex with $b_0$ connected components --- isomorphism $H_0(\Sigma; G) \cong G^{\oplus b_0}$.
\item Prove that $H_k(\Delta_m; G) = 0$ for all $k \ge 1$ and any $G$, where $\Delta_m$ is an $m$-dimensional simplex. \textbf{Hint.} Construct a contractive chain homotopy $s_k \colon C_k(\Delta_m; G) \to C_{k+1}(\Delta_m; G)$, setting $s_k\bigl([v_{i_0}, \dots, v_{i_k}]\bigr) = [v_0, v_{i_0}, \dots, v_{i_k}]$ with the removal of duplicate vertices and taking into account the sign of the resulting permutation, and check the equality $\partial s + s \partial = \operatorname{id}$ in positive dimensions.
\end{enumerate}
\end{task}

\begin{task}[nerve lemma]
Recall the nerve lemma\textup{:} if $\mathcal U = \{U_\alpha\}_{\alpha \in A}$ --- \emph{a good} open covering of $X$ \textup{(}all non-empty finite intersections of its elements are contractible\textup{)}, then the geometric realization of the nerve $|\Nerve{\mathcal U}|$ is homotopy equivalent $X$, so $H_k(X; G) \cong H_k(\Nerve{\mathcal U}; G)$ for any Abelian group $G$.
\begin{enumerate}[label=\textup{(\alph*)}]
\item The circle $S^1$ is covered by three open arcs $U_1, U_2, U_3$ with non-empty pairwise intersections and an empty triple intersection. Check that the coverage is good, describe the nerve $\Nerve{\mathcal U}$ and calculate $H_1(S^1; \Z)$ and $H_0(S^1; \Z)$.
\item Let a connected $X$ admit a good covering whose nerve is a tree, that is, a one-dimensional complex without cycles. Prove that $X$ is contractible and $H_k(X; G) = 0$ for all $k \ge 1$, and apply this to the covering of $\R^n$ by two open half-spaces $x_n < 1$ and $x_n > -1$.
\item Deduce from the nerve lemma that if a connected $X$ has a good cover of $n$ elements, then $H_k(X; G) = 0$ for all $k \ge n$ and any $G$. \textbf{Note.} The nerve of such a covering is a complex on $n$ vertices, its dimension does not exceed $n - 1$.
\end{enumerate}
\end{task}

\begin{task}[close simplicial mappings]
Two simplicial maps $f, g \colon \Sigma \to \Sigma'$ are called \emph{close} if for each simplex $\sigma \in \Sigma$ the set $f(\sigma) \cup g(\sigma)$ is a simplex of the complex $\Sigma'$.
\begin{enumerate}[label=\textup{(\alph*)}]
\item Prove that close maps induce the same homomorphisms $f_* = g_* \colon H_k(\Sigma; G) \to H_k(\Sigma'; G)$ for all $k$, and that their geometric realizations $|f|$ and $|g|$ are homotopic.
\item Show that the converse is not true\textup{:} homotopy of realizations does not imply closeness. \textbf{Hint.} Consider a three-vertex cycle $\Sigma = \Sigma' = C_3$, take $f = \operatorname{id}$ and $g$ cycling the vertices, and check that the edge image from $f$ to $g$ covers all three vertices.
\item Let $f \colon \Sigma \to \Sigma'$ and $g \colon \Sigma' \to \Sigma$ be such that $g \circ f$ is close to $\operatorname{id}_\Sigma$, and $f \circ g$ is close to $\operatorname{id}_{\Sigma'}$. Prove that $f_* \colon H_k(\Sigma; G) \to H_k(\Sigma'; G)$ is an isomorphism for all $k$.
\end{enumerate}
\end{task}

\begin{task}[stability of homology under space perturbation]
The lecture proves\textup{:} if $d_{GH}(X, Y) < r$, then for any $\varepsilon > 0$ and $\nu \ge 0$ there exist simplicial mappings
\[
\VR(Y, \varepsilon) \xrightarrow{h} \VR(X, r+\varepsilon) \xrightarrow{\iota_\nu} \VR(X, r+\varepsilon+\nu) \xrightarrow{g} \VR(Y, 2r+\varepsilon+\nu),
\]
where $\iota_\nu$ is an inclusion such that the composition $g \circ \iota_\nu \circ h$ is close to the inclusion $\iota \colon \VR(Y, \varepsilon) \hookrightarrow \VR(Y, 2r+\varepsilon+\nu)$.
\begin{enumerate}[label=\textup{(\alph*)}]
\item Based on problem \textup{3}, prove that the homomorphism $\iota_* \colon H_k\bigl(\VR(Y, \varepsilon); G\bigr) \to H_k\bigl(\VR(Y, 2r+\varepsilon+\nu); G\bigr)$ passes through $H_k\bigl(\VR(X, r+\varepsilon); G\bigr)$, that is, $\iota_* = g_* \circ (\iota_\nu)_* \circ h_*$ for all $k$.
\item Deduce from here the stability of homology\textup{:} if the class $c \in H_k\bigl(\VR(Y, \varepsilon); G\bigr)$ does not vanish when passing to the scale $2r+\varepsilon+\nu$, that is, $\iota_*(c) \neq 0$, then $h_*(c) \neq 0$. This means that a class that survives to a large scale does not disappear when moving to a space $X$ located at a distance less than $r$ from $Y$.
\item Calculate $\VR(X, r)$ and $\check C_r(X; \R^2)$ for $X = \{(0,0), (1,0), (1,1), (0,1)\}$ \textup{(}vertices of the unit square\textup{)} for $r \in (1, \sqrt{2})$ and show that both complexes are a cycle of four vertices, so $H_1 = \Z$ and $H_0 = \Z$. \textbf{Hint.} In $\R^2$, three convex sets intersect in pairs if and only if they have a common point \textup{(}Helly's theorem\textup{)}, and balls of radius $r/2$ intersect exactly when the distance between the centers is less than $r$\textup{;} also explain what happens at $r \ge \sqrt{2}$.
\end{enumerate}
\end{task}

\begin{task}[paracompactness and partitioning of unity]
A space is \emph{paracompact} if each of its open coverings can contain a locally finite open covering. Fact \textup{(}Stone, \textup{1948}\textup{)}\textup{:} every metric space is paracompact.
\begin{enumerate}[label=\textup{(\alph*)}]
\item Prove that for any open covering $\mathcal U = \{U_\alpha\}$ of the metric space $X$ there is a subdivided partition of unity\textup{:} such a locally finite family of continuous functions $\varphi_\alpha \colon X \to [0, 1]$ such that $\operatorname{supp} \varphi_\alpha \subset U_\alpha$ and $\sum_\alpha \varphi_\alpha \equiv 1$. \textbf{Hint.} Insert into $\mathcal U$ a locally finite open cover $\{W_\beta\}$, choose $\alpha(\beta)$ with $W_\beta \subset U_{\alpha(\beta)}$ and set $\psi_\beta(x) = \max\bigl\{0, 1 - \operatorname{dist}(x, X \setminus W_\beta)/\varrho_\beta\bigr\}$ for a suitable $\varrho_\beta > 0$\textup{;} then sum $\psi_\beta$ over all $\beta$ with $\alpha(\beta) = \alpha$ and normalize the resulting functions.
\item Explain how the partition of unity is used in the nerve lemma\textup{:} prove that the formula $f(x) = \sum_\alpha \varphi_\alpha(x) \, \alpha$ correctly defines a continuous mapping $f \colon X \to |\Nerve{\mathcal U}|$ into the geometric realization of the covering nerve. \textbf{Hint.} Local finiteness ensures continuity of coordinates, and the positive coordinates of the point $x$ correspond to those $\alpha$ for which $x \in U_\alpha$.
\end{enumerate}
\end{task}

\begin{task}[higher homology of connected manifolds]
Let us recall the fact used in chapter \textup{4}\textup{:} for a connected $n$-dimensional manifold $M$ without boundary
\[
H_n(M; \Z_2) = \begin{cases} \Z_2, & M \text{ compact}, \\ 0, & M \text{ non-compact}. \end{cases}
\]
\begin{enumerate}[label=\textup{(\alph*)}]
\item Prove that a compact connected $n$-dimensional manifold without boundary cannot be covered by $n$ open sets so that all non-empty intersections of these sets are contractible. \textbf{Hint.} The nerve of a covering of $n$ elements has a dimension no higher than $n-1$, which means $H_n(\Nerve{\mathcal U}; \Z_2) = 0$, whereas according to the nerve lemma $H_n(M; \Z_2) \cong H_n(\Nerve{\mathcal U}; \Z_2) = \Z_2 \neq 0$.
\item Show that compactness in \textup{(a)} is essential\textup{:} the non-compact $\R^n$ has a good cover of just two elements. Also calculate $H_n(M; \Z_2)$ for $M = S^n$, $M = \R^n$, $M = $ open unit ball, $M = $ closed unit ball, and $M = S^1 \sqcup S^1$ \textup{(}in the latter case, compare the answer with the statement of fact for a connected circle and explain why connectedness is important\textup{)}.
\end{enumerate}
\end{task}

\end{document}
